% =====================================================================
%  R. Nielsen, HVORLEDES MAN DANNER SIG EN VERDENSANSKUELSE
%  In: For Ide og Virkelighed. Et Tidsskrift, udgivet af R. Nielsen og Rud. Schmidt.  Aargang 1873 (Bind 8).
%  Kjøbenhavn: Otto B. Wroblewskys Forlag; G. S. Wibes Bogtrykkeri.  Bind 8, pp. 1-18.  Antiqua.
%
%  TRANSCRIPTION -- GENERATED; DO NOT EDIT BY HAND.  Rasmus Nielsen's eight essays in the volumes 1870-1873 are built by ONE
%  pipeline (ebook-method/for-idee-1870-73/) into a combined working file, which split7.py cuts into seven transcriptions
%  (Aandsvirkelighed and its sequel, Videre om Aandsvirkelighed, are one: Hans, 2026-10-13).
%  Rebuild from the repo root:   sh ebook-method/for-idee-1870-73/rebuild.sh
%
%  STATUS: COMPLETE (2026-10-13): the volume's title page and the whole essay, Bind 8, pp. 1-18, with its notes.
%
%  EDITION.  The only printing.  Text of record: Nasjonalbiblioteket's scan (nb.no) of the volume, ~/bibliotek/Nielsen, Rasmus/1873-for-idee-og-virkelighed.pdf (NB item 69cac3b7c05f6cf3ecc29b4b4154c2f4) (SCANS.tsv),
%  page images fetched by nbfetch.py, 300 dpi colour, NO text layer; printed page N = PDF page N + 4.  No e-book, no
%  Wikisource text, no reprint, no earlier transcription found (checked 2026-10-13: the repo, the web, the Internet Archive,
%  da.wikisource; HathiTrust and the Google Books API could not be queried that day; Google Books holds at least the 1873
%  volume, nielsen/darwinismen's source -- a second copy, not used).
%
%  ACCURACY (2026-10-13; one collation for all eight texts, 231 pp., ~58,000 words).  The DRAFT is NB's own OCR (ALTO, ABBYY
%  FineReader 12), laid over the page images (combine.py) and turned into LaTeX by draft.py; it was collated word by word
%  with an independent tesseract reading of the same images (collate/check.py, witness 'ocr'): 1556 disagreements.
%  check.py's own rules settled 78.  Candidate rules were first scored BY SCRIPT on a text already read at the image
%  (calibrate.py: Et Synspunkt for Darwinismen, same volume as Verdensanskuelse), then measured on a BLIND PILOT of 152 crops
%  (reader A, 52/52 planted controls right); 17 rules passed (NB's glyph faults: o for ø, D/O/U for Ø, i for l, li for h,
%  l for I, opening » read «, words run together or broken apart; tesseract's: aa read an/ag/å, added accents; the
%  printer's space before a mark, house rule) and settled 772; the rest were read at the image BLIND (562 crops on 94
%  sheets, four readers, 188/188 planted controls right), ~160 decided in the main session from the readers' lines and ~40
%  at 300-dpi zooms.  STRUCTURE read at the image (reader K): every footnote whole (30; collated where they stand, then moved
%  to their marks, then compared word by word with the reading: notecmp.py, 9 differences settled at the image), the
%  headings, the section numerals, the essays' last lines and signatures, the three Greek words, the volume title pages.
%  LETTERSPACING (the print's emphasis, \emph): candidate lines by letter gaps and box width, calibrated on the words the
%  word readers had marked spaced (spaced.py): 361 lines read blind (61/62 roman controls), 25 doubtful extents read at a
%  zoom.  SWEEPS after OPEN 0: words run together that both OCRs share (lower-case + capital), the unknown-word sweep (678
%  words, ~20 suspects, 14 read at the image: 4 OCR faults, 7 misprints kept), --moved, the dropped-word list, the
%  W-shortens list of the blind decisions.
%  SAMPLED PAGE CHECK (pagecheck.py, seed 202610137): one page of each essay or part, pp. 2289, 2400, 3305, 4308, 5033,
%  5208, 6514, 8011 of the working numbering (Bind x 1000 + page), read whole and blind, 1965 words: 0 word errors,
%  1 punctuation (a compound's hyphen lost at a page turn: „Ny-|Holland“) and 1 letterspacing („to be or not to be“)
%  error, both mended.  95 % upper bound for word errors: 1.5 per 1000 words.
%  What this cannot see: a misprint both OCRs normalise into a real word.
%  THE TRANSLATION AS A SECOND READING (2026-10-14): the 15 translators listed 74 doubts; ~50 were read at 300-dpi
%  zooms: 13 faults mended (4 compound hyphens set with a space, 4 specks taken for marks, a lost „I“, a colon read as
%  a semicolon, a comma read as a stop, „à“ read „å“, a space before a mark) and 9 false paragraph breaks at page
%  turns removed (left by the notes moved away; ALTO shows the next page flush); 3 more misprints kept (PRINTED AS IS).
%  The ø rule (Hans 2026-10-08): o-forms that are no word are set with ø; „stode“ (= stood) is a word and stays.
%  PRINTED AS IS: nothing recorded.
%  Not transcribed: running heads, folios, sheet signatures („For Ide og Virkelighed. 1870. I. 20“), library stamps; the
%  large initial of each essay is set as a plain letter.  Quotation marks as printed (» « in the text, sometimes a closing
%  »; „ “ in some notes).  The notes are numbered ¹), ²) per page, as printed.  Greek as printed (textalpha).
%
%  TOOLS (ebook-method/for-idee-1870-73/): nbalto.py, combine.py, alto.py, ocr.py, draft.py, rules.py, calibrate.py,
%  sheets.py + reads.py + READER-BRIEF.md, spaced.py + spaced.tsv, structsheets.py, notecmp.py, unknown.py, zoom.py,
%  pagecheck.py + PAGECHECK-BRIEF.md, patches.py (PATCHES, PRENOTES, NOTES, POST, MISPRINTS), decisions.tsv, split7.py,
%  rebuild.sh.
% =====================================================================
\documentclass[12pt,a4paper,openany]{book}
\usepackage[utf8]{inputenc}
\usepackage[T1]{fontenc}
\usepackage{libertinus}
\usepackage{libertinust1math}
\usepackage{textalpha}
\usepackage[danish]{babel}
\usepackage[protrusion=true,expansion=true]{microtype}
\usepackage{setspace}
\usepackage[margin=1.3in]{geometry}
\usepackage{fancyhdr}
\usepackage{hyperref}
\hypersetup{pdftitle={Hvorledes man danner sig en Verdensanskuelse}, pdfauthor={Rasmus Nielsen}, hidelinks}
\pagestyle{fancy}
\fancyhf{}
\fancyhead[LE,RO]{\thepage}
\fancyhead[C]{\textit{R. Nielsen: Hvorledes man danner sig en Verdensanskuelse (1873)}}
\setstretch{1.3}
\usepackage{marginnote}
\renewcommand{\thefootnote}{\arabic{footnote})}   % the print numbers each page's notes ¹), ²)
\newcommand{\opage}[1]{\marginnote{\footnotesize\textit{[#1]}}}
% the essay's heading: title (and subtitle), centred, a short rule; the print's large initial is not reproduced
\newcommand{\essayhead}[2]{\begin{center}{\large #1}\ifx&#2&\else\\[0.3em]{\normalsize #2}\fi\\[0.5em]\rule{3em}{0.4pt}\end{center}\medskip\noindent}
\newcommand{\signature}[1]{\par\smallskip\hfill #1\hspace*{2em}\par}

\begin{document}

% --- front matter: the title page of Bind 8 (PDF 1 of 1873-for-idee-og-virkelighed.pdf), read at the image (reader K, m1);
% short wavy rules are set as plain rules; the library stamps are not transcribed.
\begin{titlepage}
\begin{center}
\vspace*{3em}
{\Huge For Ide og Virkelighed.}\\[1em]
\rule{8em}{0.4pt}\\[1em]
{\large Et Tidsskrift,}\\[0.8em]
{\small udgivet}\\[0.4em]
{\footnotesize af}\\[0.6em]
{\large R. Nielsen og Rud. Schmidt.}\\[0.4em]
\rule{8em}{0.4pt}\\[0.6em]
{\large Aargang 1873.}\\[0.6em]
\rule{8em}{0.4pt}\\[4em]
Kjøbenhavn.\\[0.4em]
Otto B. Wroblewskys Forlag.\\[1.5em]
{\small G. S. Wibes Bogtrykkeri.}
\end{center}
\end{titlepage}

% ===== TEXT: pp. 1-18 =====
\essayhead{Hvorledes man danner sig en Verdensanskuelse.}{}
% --- p. 1 ---
\setcounter{footnote}{0}%
\opage{1}Hvad er en Verdensanskuelse? Noget, der nøder sig
paa En, Noget, der kan faaes for godt Kjøb, og Noget,
som maa kjøbes til den dyreste Pris, alt som man tager
det. Alle Mennesker have Verdensanskuelse, hver har
sin; de kunne ligne hinanden saa paafaldende, at man
skulde sværge paa, de vare vævede i eet Mønster, men de
kunne ogsaa være høist forskjellige. Individer kunne være
uenige om mange Ting og dog have samme Verdensanskuelse;
de kunne ogsaa enes om mange Ting, og dog
være uenige i Verdensanskuelse. Det er for Samfundet
vigtigt at vide, hvormegen Enighed der i denne Henseende
er nødvendig, naar den borgerlige Orden skal opretholdes;
thi vel maa ethvert Menneske have Ret til at
see Verden og Livet med sit Syn, men gaaer Uenigheden
mellem Individerne til Yderligheder, saa er Samfundet
i Fare. Der er virkelig Verdensanskuelser, der
ikke kunne forliges; Tyven f. Ex. har een Verdensanskuelse,
Politiet har en anden, saadanne to lade sig ikke »mediere«.

Dersom det ikke var for Verdensanskuelsens Skyld,
vilde det være en sand Besparelse at afskaffe Kultusministeriet;
thi hvad er vel dette Ministeriums Øiemed i
% --- p. 2 ---
\setcounter{footnote}{0}%
\opage{2}Grunden Andet end at drage Omsorg for, at en vis almindelig,
Samfundet betryggende og desaarsag for alle
Statsborgere, Smaa og Store, høist nødvendig Verdensanskuelse
kan blive opretholdt, udviklet og paa indtrængende
Maade komme til hele Folkets Bevidsthed. Men
for at Besparelsen skal kunne indføres uden Fare for
Samfundet, maa der først gjøres en ny og stor Opdagelse.
Der maa absolut opfindes et Arkanum, et Tryllemiddel,
ved hvis Anvendelse man lige med Et bliver i Stand til
at indplante i alle Sjæle, i alle Menneskers Sind og Hjerte
en lige klar og levende Sands for Retfærdighed, Sandhed
og Billighed, et fælles ubedrageligt Instinkt, som siger
Enhver, at Ingen kan have Rettigheder uden at have
Pligter eller borgerlig Frihed uden at underkaste sig Lovene,
en varm og stærk og inderlig Følelse af, at vi alle
ere Brødre og Søstre og skulle i al Oprigtighed elske
hverandre, saa at Ingen seer paa sit Eget uden tillige
at see paa det fælles Bedste: ja, kunde Nogen blot
opfinde et saadant Arkanum, da skulde Slægterne juble
og Tidens Aander hilse den nye Messias, thi da vare de
haardeste Samfundsknuder løste. Hvilke Besparelser af
Penge, Bekymring og aandeligt Arbeide! Saa kunde man
da rigtignok afskaffe baade Kultusministeriet og Skoletvangen
og Politiet og Præstelønningen og det kongelige
Theater og Hærloven; man kunde i Sandhed indføre
gjennemgribende Besparelser. Og hvilke Fremskridt i
Friheden! Vilde En lære Noget, var det godt; vilde han
Intet lære, var det ogsaa godt; man kunde nu stole paa
hans Instinkt og Følelse. Vilde En for Salighedens Skyld
have sig en Religion, det være en brahminsk, muhamedansk,
jødisk eller kristelig, naar han blot ikke vilde paanøde
% --- p. 3 ---
\setcounter{footnote}{0}%
\opage{3}Andre sin Tro, saa var det godt; vilde han ingen Religion
have, men ansee Saligheden for en Luxusartikel, var
det ogsaa godt; man kunde jo stole trygt paa hans Instinkt
og Følelse. En saadan fuldkommen Frihed med
Hensyn paa den religiøse Verdensanskuelse vilde naturligvis
medføre en lignende Frihed med Hensyn paa den
sædelige. Ligesom Enhver bør have Lov til at klæde sig
efter sin Smag, kun at han ikke gaaer ganske nøgen, saaledes
skulde da ogsaa Enhver have Lov til at leve privat
paa sin egen Vis; Instinktet og Følelsen vilde jo borge
for, at det skete i god Forstaaelse med de Medlevende.
Vilde En være Mormon og have sytten Koner, var det
godt; vilde han kun have en eneste, var det godt; vilde
han ingen have, men nøies med uvisse Indtægter, var
det ogsaa godt. Saa vilde Friheden trives og »Selvstyret«
blomstre, --- naar vi blot havde Tryllemidlet!

Da et saadant Tryllemiddel endnu, saa vidt vides, ikke er
opfundet, saa er Staten nødsaget til at betjene sig af de Opdragelsesmidler,
som, i Overensstemmelse med Befolkningens
Traditioner, Tro og Sæder, ifølge den bestaaende Forfatning
vise sig at være de virksomste, og det er unegtelig
dem, som findes i Kirken og Skolen. Det bør udtrykkelig
fremhæves, at Staten, skjøndt den fra sit Synspunkt fornemmelig
maa betragte disse Institutioner som Midler
til Opretholdelse af borgerlig sædelig Orden, dog derfor
ingenlunde miskjender, at de have Værd og Gyldighed
som Maal i sig selv; men det bliver nødvendigt at
gjøre en skarp Adskillelse mellem det, der vedkommer
Kirken som Kirke, Skolen som Skole og Staten som Stat,
at det deraf kan sees, hvormeget der er Gjenstand for
Tvang og hvad der, ifølge Sagens Natur, tilhører Friheden.

% --- p. 4 ---
\setcounter{footnote}{0}%
\opage{4}For Kirken som Menighed er Religionen eneste og
høieste Formaal, og Troen er fri, Ordets Forkyndelse fri;
alle Menigheder ere fødte Frimenigheder, der gives slet
ingen Undtagelse. Ikke desto mindre maa Staten betragte
Religionen som et Opdragelsesmiddel og vurdere den
religiøse Verdensanskuelse fra et humant, statsborgerligt
Synspunkt. Den maa for Opdragelsens Skyld drage Omsorg
for, at der sikkres Religionslæreren en saa fri og
selvstændig Stilling over for Tilhørerne, at Sandhedsordet
kan blive hørt og Ordets Forkynder ikke synke ned
til at blive Klient af det Flertal, han er sat til at opdrage.
Religionslæreren maa --- for Opdragelsens Skyld ---
staae Staten til Regnskab, og hans Gjerning bedømmes, ikke
efter den særegne Troesbekjendelse, hvorpaa han staaer,
men efter den Dygtighed, hvormed han formaaer at vække
i Sjælene den levende Sandhedserkjendelse, den dybe Ærefrygt
for Ret og Pligt, som er uundværlig for gode Statsborgere,
og gjøre det praktisk paa en saadan Maade, at
Læren bliver til Overbevisning, og Overbevisningen til
Kjød og Blod i Menighedens Lemmer. Staten bør i
denne Henseende føre et strengt konsekvent Overtilsyn
med alle Religionslærere i alle Menigheder af alle Troesbekjendelser:
der gives --- for Opdragelsens Skyld ---
ingen Undtagelser.

For Skolen er Undervisningen en Hovedsag. Enten
der nu nærmest lægges an paa Vækkelse eller paa Kundskabsmeddelelse,
har Skolen som Skole sit eget Maal og
i det selvegne Maal sin Frihed. Alle Skoler, Almueskolerne
lige saa vel som de lærde, ere i denne Henseende
at betragte som Friskoler; ingen udenfor staaende
Myndighed bør gribe vilkaarlig ind i Skolens Plan: jo
% --- p. 5 ---
\setcounter{footnote}{0}%
\opage{5}friere det kan gaae, desto bedre. Men Skolerne ere tillige
Opdragelsesanstalter; hvorledes Undervisningen end drives,
maa der, naar den da ikke er ganske karakterløs, naturlig
danne sig en Forestilling om det fælles Menneskelige, et
Syn paa Livet, en Slægts Verdensanskuelse. Paa dette
Punkt er Staten opmærksom. Hvad der i det Enkelte er
Undervisningsgjenstand, hvormeget eller hvor lidt der
læses, er ikke Hovedspørgsmaalet; det for Staten afgjørende
Spørgsmaal er, hvorledes det gaaer med Opdragelsen.
Bibringes der de Unge en sund og sand og fyldig Erkjendelse
af, hvad en dygtig Statsborger er sig selv og
sit Fædreland skyldig, saa er Alting godt; skeer det ikke,
saa er Alting slet, om saa alle Eleverne gik sprænglærde
ud af Skolen. Skolens Lærere og Bestyrere maae i dette
Punkt staae Staten til Regnskab; at beraabe sig paa Friheden
og give Forsikkringer om Instinkt og Følelse,
hjælper ikke: Staten kan ikke nøies med det Uvisse; den
maa have fuld Vished, den maa --- for Opdragelsens
Skyld --- føre et virksomt Overtilsyn med alle Skoler
uden Undtagelse.

Skulde det nu synes, som om vi indrømmede Staten
et Formynderskab, der maatte være trykkende for Aandsfriheden
i Kirke og Skole, da bede vi bemærket, at det,
hvorpaa Formynderskabet gaaer ud, er for Frihedens Skyld
indskrænket til det Allernødvendigste, det Alleralmindeligste,
til Noget, som kan naaes ad meget forskjellige
Veie. Det maa dog Alle tilstaae, at en Religion, som,
hvormeget der end prækes, ikke engang formaaer at vække
den borgerlige Samvittighed, er en maadelig Religion, og
en Skole, som trods Lektier, Overhøringer og Examiner,
ikke er i Stand til at indprænte og befæste i de Unges
% --- p. 6 ---
\setcounter{footnote}{0}%
\opage{6}Sind de første nødvendige Grundsætninger om Ret og
Pligt, en maadelig Skole. I ældre Dage vilde Staten
for Meget; den vilde gjøre Skoledrengene til Dogmatikere
og Undersaatterne til Kristne, det mislykkedes. I vore
Dage, da Staten har begyndt at indskrænke sine Fordringer
til det Nødvendigste, er der Haab om, at den
netop vil naae, hvad den tilsigter.

Den Verdensanskuelse, som Staten for enhver Pris
vil have opretholdt, kan udtales i faa Ord: det er Menneskets
Bestemmelse at leve i et retlig ordnet Samfund,
adlyde Lovene og udfylde sin Plads ved en samvittighedsfuld
Opfyldelse af sine Borgerpligter; den er i sin Almindelighed
saa farveløs, at den ikke tager sig ud paa
Papiret, men i Virkeligheden saa vigtig, at al Begeistring
for Folkelighed, Frihed og Fædreland uden den bliver til
Luftblærer og blaa Dunster. Den sociale Verdensanskuelse
har adskillige Fjender, den maa bekæmpe. Den kan f. Ex.
ikke bestaae med et paveligt Hierarki, thi Hierarkiet sætter
sig ud over de statsborgerlige Love; den kan ikke bestaae
med en asketisk Livsretning, der i Armod og Forsagelse
foragter det Timelige; den kan ikke bestaae med en selvisk
Raahed, der trodser Lovene og stormer mod Rettens
Skranker for at tage sig selv til Rette; den kan ikke
bestaae med et sneversynet, lavsindet Partivæsen, der kun
vil mætte en enkelt Stands Egoisme, være sig en Parfume-Egoisme
eller en Tran-Egoisme; den kan endelig ikke bestaae
med den bundløse Alsidighed, som under Mærket:
\emph{finere Dannelse}, lever af at lugte til Alt, snakke om
Alt, være færdig med Alt og tvivle om Alt, kun ikke om
Nydelsens, Adspredelsens og Opinionens intelligente Guddommelighed.
Naar disse Modstandere fraregnes, vil den
% --- p. 7 ---
\setcounter{footnote}{0}%
\opage{7}officielle, statsborgerlige Verdensanskuelse kunne forliges
med en hvilkensomhelst personlig sædelig Overbevisning.

Den personlige Verdensanskuelse, der skal bære den
statsborgerlige, faaer fornemmelig Lys og Varme fra de
to Brændpunkter, Guddommen og Sjælen. Er der en
Gud, og hvad er Sjælen? Dette Tvillingspørgsmaal paatrænger
sig Enhver, der begynder at tænke over Tilværelsen.
Kunne slige Spørgsmaal besvares? Igjen et
Spørgsmaal. Hvilken Vei skal man egenlig gaae for at
faae dem besvarede? Igjen et Spørgsmaal. Saa vaagner
Tvivlen ved de mange Spørgsmaal. Vor Tid er en Tvivlenes
Tid. At tvivle om Alt er, mener man, et Tegn paa
overlegen Forstand, det er Visdommens Begyndelse. Og
dog har Tidsalderen med den travle Kritik og den frie
Tanke kun liden Glæde af al sin Tvivlen; den er forstemt,
thi Tvivlen tærer; den føler sig mat, mismodig og lider
af aandelig Hunger. At tvivle for ret at prøve kan saavel
i praktisk som i theoretisk Forstand være sundt, naar
det der skal prøves, er af den Natur, at Prøvelsen kan
ende Tvivlen. Men at tvivle personlig om sit eget Selv,
om Livets Betydning og Verdens Ophav, nære Tvivlen,
ruge over Tvivlen, døse i Tvivlen, er usundt; en saadan
Tvivlen er Sygdom. Nutiden lider af Tvivlesyge; Tankernes
Puls er urolig, Verden har Krampe, Slægten har
Feber; det er Selvmodsigelsens Orme, der gnave paa Hjertets
Rødder. Under stærke Kuldegysninger faae de urimeligste
Forestillinger Lyd, og den ene Modsigelse jager den
anden. Det moderne Menneske veed bogstavelig ikke,
hvad det skal antage sig selv for, et Dyr, et Stykke Gud
eller begge Dele. Paa den ene Side føres der et Sprog
i Kritik og Filosofi, som havde den menneskelige Viden
% --- p. 8 ---
\setcounter{footnote}{0}%
\opage{8}ingen medfødt, uoverstigelig Grændse; Verdensgaaderne
løses, den gamle Mulighedernes Nat maa ud med Hemmelighederne:
Tilværelsen har intet Mysterium, det Absolute
selv ligger aabent, gjennemsigtigt for Tænkerens
Blik. Det kan man da forstaae, dersom det virkelig er
sandt, at Aanden først bliver Aand i Mennesket; det
kan man forstaae, naar det gjælder følgestrængt, at Guddommen,
dette Alnaturens Princip med den blinde Visdom,
hvis Lys er Mørke, er et Bevidsthedsprincip, der
stræber op mod Selvbevidsthed, men ikke kan blive sig
selv bevidst, før det ved Materiens Hjælp faaer bygget
sig en menneskelig Hjerne. Fra dette Synspunkt er da
Mennesket den selvbevidste Guddom, og den selvbevidste
Guddom maa i Videnskaben --- efterhaanden, forstaaer
sig, gjennem mange Anstrængelser og kun i Tilnærmelse!
--- vel kunne gjennemskue sit eget ubevidste Værk. Nu
skulde man da antage, at det i Videns Navn var Alvor
med Bevidstheden, at den filosofiske Selvbevidsthed, det
absolute Indholds absolute Form, maatte være det i sig
selv Høieste, Herligste, Allerfuldkomneste; men Nei! det
er ei saaledes ment. Mennesket er, det maa for Alting
ikke glemmes, et Dyr, et videnskabeligt Dyr. Kun som
Dyr har det Bevidsthed, kun i Kraft af Bevidsthed har
det Selvbevidsthed: Dyrets Selvbevidsthed er lige saa forgængelig
som Dyret selv. Man seer, hvorledes Belysningen
skifter. \emph{Først skulde det Bevidste være høiere end
det Ubevidste}; det Absolute maatte gaae den lange
Omvei, inden det fik Bevidstheden frem, thi det søgte
Viden, og kun i Bevidstheden var der en Viden. Men
saa kommer det pludselig bag paa os: Bevidstheden har
kun en indskrænket, af Hjernesvingninger afhængig Viden;
% --- p. 9 ---
\setcounter{footnote}{0}%
\opage{9}kommer der Uorden i Hjernesvingningerne --- hvad desværre
let kan skee --- saa kommer der ogsaa Uorden i
Viden; fik en Tænker Vand i Hovedet, maatte det vidende
Absolute jo ogsaa faae Vand i Hovedet; det kan den
evig guddommelige Viden ikke være tjent med. Hvad
skal da Viden gjøre? Emancipere sig fra alle Hoveder,
alle Hjerner og Hjernesvingninger, altsaa fra Bevidstheden!
det var en snild Vending. \emph{Nu er det Ubevidste
høiere end det Bevidste}. Det Absolute har da ingen
Bryderier med Hjernesvingninger; det kan undvære Hjernerne,
thi det kan undvære al Bevidsthed. Om alle Tænkere
kom i Daarekisten, det Absolute kom dog ikke i
Daarekisten, thi det kan undvære alle Tænkere. Den
fuldkomne Viden er en Viden i det Ubevidste, en salig
Viden. Hvor trøsterigt ikke blot for det Absolute, men
ogsaa for Dyret med Hjernen. Nu kan det vidende Dyr
døe, naar det vil; det gaaer rigtignok ud over Hjernen,
thi Hjernen er jo Bevidstheden; men det Vidende i Hjernen
tager ingen Skade, tvertimod! den hjerneløse Viden
er en salig Viden, salig i det Ubevidste. At her er to
uforligelige Verdensanskuelser, vil Ingen kunne nægte, og
dog maae de forliges og smelte sammen til een, thi
ellers kommer Dualismen.

Der hører Konstfærdighed til, naar man skal afsætte
en firkantet Cirkel; og der hører Skarpsindighed til og
et fint Kjendskab til Filosofiens Historie, naar man skal
forlige to saa uforligelige Anskuelser som ovenstaaende.

Gaaer man ud fra den Antagelse, at Hjerne og
Bevidsthed hos Mennesket maae staae og falde med hinanden,
da har man kun Valg mellem Et af To. Enten
indrømme, at der gives en guddommelig Bevidsthed,
% --- p. 10 ---
\setcounter{footnote}{0}%
\opage{10}der er uendelig ophøiet over den menneskelige, eftersom
der gives en guddommelig Viden, der er uendelig ophøiet
over den menneskelige; i dette Tilfælde er da den menneskelige,
af Hjernen betingede Videns Grændse en Gang
for alle erkjendt, og der er et Tilværelsens Mysterium.
Eller aabenlyst erklære, at der i Tingenes Væsen slet
ingen Tanker er, eftersom der i Tingene selv ingen Bevidsthed
er, og at ubevidste Tanker ere umulige. --- Der
gives i Sjælen dunkle Tanker, fordi der gives en dunkel
Bevidsthed og en dunkel Tænken. Tanken er dunkel,
naar Bevidstheden er dunkel; men ingen Tanke, den være
klar eller dunkel, er mulig uden en tilsvarende Bevidsthed.
For at tænke os det Ubevidste, maae vi lige saa
nødvendig bruge vor Bevidsthed som for at tænke os det
Bevidste. Vi sige om en Sten, at den er ubevidst, men
vi sige det i Kraft af vor Bevidsthed. Den, der aldrig
har seet Lyset, har ingen Forestilling om Mørket, og
den, der aldrig er blevet sig Noget bevidst, veed ikke af,
at der er noget Ubevidst eller, hvad det Ubevidste er:
han har slet ingen Viden.

Hvis en Fyrste eller en rig Englænder en Gang fik
i Sinde at udsætte et Par Millioner Pund Sterling som
Præmie for Besvarelsen af det Spørgsmaal: \emph{Hvorledes
bærer den menneskelige Bevidsthed sig ad med
at tænke en fuldkommen Viden uden al Bevidsthed}?
saa blev det da ikke en Tænker af almindelig Organisation,
det maatte blive en ægte Absolut, der vandt
Præmien. Det er en let Sag at underlægge Naturkræfterne
en ubevidst Villie og tale om Naturens ubevidste
Øiemed, naar derved ikke skal tænkes paa Andet end en
billedlig Overførelse af det, der stammer fra Bevidstheden,
paa det, som ingen Bevidsthed har. Kalder man uden
% --- p. 11 ---
\setcounter{footnote}{0}%
\opage{11}videre al Stræben Villie, saa er der Villie overalt, Villie
i Stenen til at falde, Villie i Maanen til at nærme sig
Jorden o. s. v. Man kunde for den Sags Skyld gjerne
sige, at der i Lokomotivet paa en Jernbane er en ubevidst
Villie til at trække Vaggonerne, forenet med en ubevidst
Forestilling om den næste Station, der skal naaes.
Men det er kun Talemaader. Netop den Kjendsgjerning,
at Naturen virker som efter Øiemed, uagtet den ikke har
noget Øiemed, at den saa forunderlig minder om en Alvidende,
skjøndt den slet ingen Viden har, er den store
Gaade. Naar den menneskelige Viden, for at kunne
rumme det Absolute, absolut maa af med Bevidstheden
og dog ikke kan blive Bevidstheden kvit, saa er den ilde
faren; Tvivlene reise sig, Modsigelsen piner, og Kvalen
kan blive tantalisk.

Enhver, som for den personlige Overbevisnings Skyld
føler Trang til at danne sig en Verdensanskuelse, see vel
til, at han ikke blot har Tankerne frie, men ogsaa Sindet
frit. Uden frie Tanker vil han paa mange Punkter hilde
sig i det hidtil Vedtagne, i indgroede Forudsætninger,
som ikke blive prøvede, og idet han dog lægger Adskilligt
til af sit Eget, vil Gjennemarbeidelsen af det Hele
blive saa ufuldkommen, at Sammenhængen mangler.
Men selv om Tænkningsarbeidet er udført med en saadan
Følgestrænghed, at en nøieregnende Kritik vanskelig skal
kunne paavise væsenlige Uoverensstemmelser mellem Bygningens
enkelte Dele, er det endnu ikke noget fuldgyldigt
Bevis for dens Sandhed som Verdensanskuelse. En Verdensanskuelse
kan i formel Henseende have videnskabelig
Konsekvens og god Sammenhæng og i Realiteten endda
være falsk, grundfalsk, ja urimelig indtil det Vanvittige,
% --- p. 12 ---
\setcounter{footnote}{0}%
\opage{12}naar Tænkerens Sind er hildet og ufrit. Vi kunne herpaa
anføre et storartet Exempel fra den nyeste Filosofi i
Tydskland.

»Det Ubevidstes Filosofi«;\footnote[1]{Philosophie des Unbewussten. Versuch einer Weltanschauung von E. v. Hartmann, Dr. phil. Berlin 1869.} en Bog, der allerede, og
med Grund, har vakt en betydelig Opmærksomhed i den
tydske Literatur, er mere end en velskreven Bog; det
er et i fine Træk og med skarpsindig Analyse af Naturlærens,
navnlig Fysiologiens, allervanskeligste Problemer
gjennemført Aandsarbeide, som her foreligger. Paa en
Bedømmelse af Værkets videnskabelige Anlæg og eiendommelige
Udgangspunkt ville vi i nærværende Sammenhæng
ikke indlade os; kun af den Verdensanskuelse, som
heri udtales, anføre vi nogle Træk til en Prøve. --- Lige
fra Schellings første ungdommelige Optræden med sin
intellektuelle Anskuelse af det Absolute indtil den seneste
Tid har den spekulative Filosofi i tvetydig Forstaaelse
med Naturlæren arbeidet paa at klare sig det absolute
Væsen. Det har imidlertid været det intelligente, medfilosoferende
Publikum vanskeligt at følge med og komme
efter, hvordan det egenlig hang sammen med dette
Absolute. Snart blev det fremstillet som en logisk alvidende
Guddom: Gud i Logiken, før han skabte Verden,
snart igjen som et uoplukket, dunkelt Princip, Altilværelsens
iboende Grund og Aarsag. Nu er da endelig
Tvetydigheden hævet; Sandhedsgudindens Aasyn sees fuldstændig
afsløret. Det Ubevidstes Filosofi, der ikke ensidig
holder sig til Spekulationen, men nøiagtig reviderer Erfaringens
Kjendsgjerninger og tager kraftige Tag i Mate%
% --- p. 13 ---
\setcounter{footnote}{0}%
\opage{13}rialismen, har udgrundet Hemmeligheden. Det store
Spøgelse, som saa længe har huseret i spekulative immanente
Hjerner, er manet frem, det Absolute viser sig ved
høilys Dag og det Absolute er? --- det rene Intet, det
evigt salige »gar Nichts«.

Verdensgaaden er, at en ubevidst Villen fra Begyndelsen
af har drevet en ligeledes ubevidst Forestillen fra
Ikke-Væren ind i Væren, thi, hedder det S. 607, »i sig
havde Forestillingen ingen Drift og ingen Interesse af at
træde fra Ikke-Væren over i Væren, følgelig var før Indtræden
af Villen ingen Forestillen aktuel, følgelig var der
før Verdens Tilblivelse hverken Villen eller Forestillen
d. e. »gar Nichts«. Saalænge Villen vedvarer, saalænge
vil Processen og dens Tilsynekommen i Bevidstheden,
Verden, vedvare; naar der altsaa en Gang ikke skal være
nogen Verden mere, da tør der heller ikke mere være
nogen Villen eller nogen Forestillen (da den ubevidste
Forestilling altid kun forsaavidt bliver aktuel, som Villiens
Interesse fordrer den) d. e. der vil igjen være »gar
\emph{Nichts}.« Hermed stemmer da ogsaa følgende Beskrivelse
af Guddommen (S. 606): »Det Al-Eneste Ubevidste
er alvidende og alviist, og kan følgelig ikke blive klogere
end det allerede er; det har, som ogsaa Aristoteles siger,
ingen Hukommelse, altsaa kan det af de Erfaringer, det
ellers gjorde i Verden, Intet lære. Følgelig er det, naar
Verden en Gang har ophørt at være, akkurat forblevet
det samme, som det var før Verdens Skabelse; saa saligt,
som det forhen var, er det nu igjen, hverken mere eller
mindre; Verdensprocessen kan paa ingen Maade forhjælpe
det til større Salighed, end det forhen havde\dots{}.«

Er Guddommen salig i Intet, da følger det af sig
% --- p. 14 ---
\setcounter{footnote}{0}%
\opage{14}selv, at Sjælene ogsaa maae blive salige i Intet. Haabet
om en Sjælens individuel Vedvaren efter Døden erklæres
ligefrem for en Illusion. »Dermed er de kristelige Forjættelsers
Hovednerve overskaaret, thi Mennesket er det
i Grunden dog kun om hans eget kjære Jeg at gjøre;
hvad hjælper mig den største tilkommende Salighed, naar
jeg ikke kan fornemme og nyde den!«

Man skjønner heraf, hvad Verden er, hvad Livet er
værd og hvad al menneskelig Digten og Tragten har at
betyde; det er Skuffelse, Indbildning, Illusion Altsammen.
»Det Ubevidstes Filosofi« angiver os tre Bevidsthedstrin
svarende til tre Stadier paa Illusionernes Vei. Illusionens
første Stadium er betegnet ved, at Individet bilder sig
ind at kunne ifølge en given Verdenstilstand opnaae Lykken
i det jordiske Liv. At det er et Bedrag, følger deraf,
at al Lyst kun er en Betingelse for Ulyst, al Nydelse for
Lidelse, al Glæde for Sorg og Smerte. Naar Menneskene
blot ville gjøre sig det klart, at Glæden forholder sig til
den overstandne Sorg som Smertens Ophør til Pinebænkens
Kvaler, ville de snart blive kurerede for den om
sig gribende Nydelsessyge. »Hvad man nutildags kalder
Pauperismens Spøgelse, er intet Andet end den i Masserne
opdukkende Bevidsthed om, at Kampen med Nøden,
Sorgerne og deres Lindring ligger ganske paa den negative
(Smerte-) Side af Fornemmelsens Nulpunkt, medens
tidligere, da Massens Armod var ti Gange større, denne
Bevidsthed manglede, og Folket bar sin Fattigdom som
af Guds Naade!« Paa Illusionens andet Trin bilder Individet
sig ind at skulle naae Lykken i et hinsidigt Liv
efter Døden. Det er det kristelige Bedrag. »Jesu Verdensanskuelse
var for naiv og barnlig til at adskille Sjæl
% --- p. 15 ---
\setcounter{footnote}{0}%
\opage{15}og Legeme«, hvorfor ogsaa »Kjødets Opstandelse« er med
i den tredie Troesartikel. At det er en Illusion, forstaaes
af sig selv; ja, der findes endogsaa Steder i Skriften, som
antyde, at al Bevidsthed lige som al Tid ophører i Evigheden.
Illusionens tredie og sidste Stadium er det, hvorpaa
vi kjøre som paa Jernbane. Indbildningen er, at
Lykken ligger i en stor, politisk social Fremtid. Man
venter sig Frelsen af en forbedret Samfundsorden, af
Fremskridt i Industri og Videnskab: et uhyre Bedrag!
Enden bliver, at den gamle, udlevede Menneskeslægt er
træt til Døden, træt af Arbeiderne og af Nydelserne,
uden Arvinger og uden Kjærlighed, træt af Alt. Da vil
den komme hin ophøiede Melankoli, som man sædvanlig
finder hos Genier og aandelig høitstaaende Oldinge; da
vil den komme, den rene Himmelklarhed, hin kjølig guddommelige
Ro, som vifter i Spinozas Ethik, hvor Lidenskaberne
sluges af Fornuftens Afgrund, fordi de klart og
tydelig ere fattede i Ideer. Da er Illusionen død, Haabet
udbrændt, thi hvorpaa skulde man endnu haabe? Da
erkjender den dødtrætte Menneskehed kun Et, der er
saligt, Hvile fra Sorg og Savn, Smerteløshed i det rene
Intet.

Er det nu en trøstesløs, fortvivlet Verdensanskuelse?
Ingenlunde! at sige, naar Sindet er frit og uhildet som
Tanken. Overfor en objektiv videnskabelig Theori maa
Sindet være frit. At Sindet er frit, kan betyde, at Tænkeren
ikke gjør Theorien til personlig Overbevisning, men
alene har Interesse af at undersøge, hvorpaa den er bygget,
og hvorvidt den i alle Dele hænger sammen. Det
kan ogsaa betyde, at han gjør Theorien til sin personlige
Verdensanskuelse, naar han kun veed med sig selv, at
% --- p. 16 ---
\setcounter{footnote}{0}%
\opage{16}han gjør det i Kraft af et Valg. Heri er intet Trøstesløst.
Thi hvorfor \emph{vælger} han Anskuelsen, uden fordi
han føler, at det netop er en Anskuelse for ham, en
Verdensanskuelse, der dybt og væsenlig tilfredsstiller hans
indre Trang, en Anskuelse, hvori hans matte Sjæl finder
Hvile. Hvis Nogen derimod, skjøndt han i sit Inderste
føler Trang til en levende Sandhed, en sand Erkjendelse
af det evige Liv, lader sig overvælde af en saadan Dødens
og Tilintetgjørelsens Visdom, fordi han er imponeret
af den formentlig objektive »Wahrheit«, da er han i
Sandhed trøstesløs. Men det Trøstesløse ligger ikke i
Sandheden; nei det ligger i hans eget Selvbedrag. Hans
Hjerte er sygt, hans Sind er hildet; han bilder sig ind
at staae Ansigt til Ansigt med den forfærdeligste af alle
Magter, Sandheden selv, og saa er det slet ikke Sandheden!
Theoriens objektive Gespenst er netop Illusionen,
ja Illusionernes Illusion, det afmægtigste Væsen af Alt,
Skyggen af et Intet: den objektive fra guddommelig og
menneskelig Bevidsthed skilte »Wahrheit« er --- »gar
Nichts«.

At danne sig en Verdensanskuelse er ingenlunde det
Samme som at danne sig en rigtig Forestilling, et saavidt
mulig sikkert og tydeligt Begreb om den udvortes
virkelige Verden. Erkjendelsen af Tingenes Love, Historiens
Gang og Verdens Sammenhæng kan hos en Iagttagende
være meget udviklet, og de Forudsætninger, hvorpaa
en personlig Overbevisning skal bygges, endda være
yderst dunkle, mangelfulde, forvirrede. For at erkjende
det Sande, maae vi erkjende det Virkelige og lære at
betragte Kjendsgjerningerne, være sig Naturens eller Historiens,
i deres eget Lys; det er jo vist nok. En Verdensanskuelse
saa lyssky, at den ikke tør see Virkelig%
% --- p. 17 ---
\setcounter{footnote}{0}%
\opage{17}heden lige i Øinene, saa forkjælet subjektiv, at den lister
af Veien for Videnskaben, fordi den er angst for at møde
det Objektive, saa selvisk indbildsk, at den ordenlig bærer
Nag til Naturlære, Filosofi og Kritik, fordi den har en
Mistanke om, at disse Udluftende let kunde faae i Sinde
at gjøre Brug af Feiekosten og rydde lidt op i dens
Ammestue, en saadan Verdensanskuelse er naturligvis
uden Værd for dem, der søge Lyset og elske Sandheden.
Men Virkelighedserkjendelsen er kun Sagens Indledning,
ikke Sagen. Til en personlig Verdensanskuelse hører
nødvendig et personligt Selv. Den, der glemmer sit eget
Selv, forglemmer det Allervigtigste; bedaaret af det Objektive,
har han allerede ved Udgangspunktet begyndt med
at bedrage sig selv og hilde sig i Illusionen.

Nu er det vistnok saa, at en Verdensanskuelse ikke
kan være vilkaarlig, at det menneskelige Selv ikke kan
skabe sig, som det vil, fordi det ikke kan omskabe den
virkelige Verden og forandre Naturens Orden; der er en
Livets Grundlov, der gjælder for alt Levende uden Undtagelse;
kun een Aabenbarelse af det personlige Selv kan
i evig Forstand være den fuldkomment sande. Men det
er ogsaa vist, at hvert enkelt Menneskes Erkjendelse af
Selvet er betinget af dets Selverkjendelse, og denne Selverkjendelse
af dets Selvudvikling. Da det nu er vitterligt,
at Individerne i Samfundet befinde sig paa meget
forskjellige Trin af Selvudvikling og Selverkjendelse, og
da det er umuligt, at en personlig overbevisende Verdensanskuelse
kan finde Indgang hos Nogen uden at gaae
ad Selverkjendelsens Vei: saa følger heraf, at de mange
indbyrdes stridige Verdensanskuelser, Materialisme, Nihilisme,
Atheisme, Pantheisme, Rationalisme, Jødedom,
% --- p. 18 ---
\setcounter{footnote}{0}%
\opage{18}Kristendom o. s. v., alle maae have fuld Ret til at bestaae
og udtale sig i et paa Frihedens Grundlag ordnet Samfund.

Hvad har saa Staten i Forhold til de mange personlige
Verdensanskuelser at gjøre? Opgive en gammel
Fordom; forsage en tusindaarig, indgroet Illusion, og
forsage den med samme Oprigtighed, som en troende
Baptizand forsager Djævelen og alt hans Væsen og alle
hans Gjerninger, --- den gamle sjælefordærvende Illusion,
at en ensartet, fuldt personlig Verdensanskuelse skulde
være Noget, der ved objektive Statslove kunde paatvinges
Alle eller dog de Fleste; at der skulde gives en objektiv
religiøs Videnskab, en Troesvidenskab, der kunde komme
Statsmyndigheden til Hjælp med at faae den objektive
Sandhed almindelig indført; at Naturlære, Filosofi og
Kritik kun skulde ansees for midlertidig tolererede, »hedensk
opblæste Videnskaber«, som først kunne faae Syndernes
Forladelse, naar de igjen bekvemme sig til at
knæle for Theologien og lade sig døbe paa Dogmatiken.
Ved grundig at forsage denne gamle Illusion og kraftig
sørge for Ungdommens sædelig statsborgerlige Opdragelse,
ved Overtilsyn med den hele saavel religiøse som humane
Undervisning vil Staten naae sit Maal og føle sig betrygget.

De mange Verdensanskuelser ere personlige Livsmagter;
de stride med Livets Vaaben for Livets Ret.
Alt som Striden gaaer, vil det vise sig, hvilken iblandt
dem der har Kræfterne, fordi den har Betingelserne.
Det skal vise sig, at de ikke ere lige stærke, fordi de
ikke ere lige sande; det skal i Kjendsgjerninger vise
sig, at der kun er een, der er den evig stærke, fordi
der kun er een, der er den evig sande, en eneste, hin
Evangeliets Perle, der maa kjøbes til den dyreste Pris.
\signature{R. \emph{Nielsen}.}

\end{document}
